Let \(V=\operatorname{span}(\mathcal F)\), viewed as a linear subspace of \(\mathbb R^S\) with the supremum norm, and put \(r=\dim(V)\le d\). Since every member of \(\mathcal F\) takes values in \([0,1]\), \(\mathcal F\) lies in the unit ball of \(V\). For \(0<v<1\), a maximal \(v\)-separated subset of \(\mathcal F\) is a \(v\)-net. The open supremum-norm balls of radius \(v/2\) centered at its elements are disjoint and are contained in the ball of radius \(1+v/2\). Comparing \(r\)-dimensional Lebesgue volumes on \(V\) gives
\[
N_{\mathcal F}(v)\le \left(1+\frac{2}{v}\right)^r\le \left(1+\frac{2}{v}\right)^d.
\]

We next control the dependence of \(H_{g,f}\) on its two arguments. Suppose \(\lVert g-g'\rVert_\infty\le v\), and define \(g_t=(1-t)g+tg'\). Differentiation of the finite Gibbs formula yields, for every \(x\in\mathcal X\) and \(a\in\mathcal A\),
\[
\frac{d}{dt}\pi_{g_t}(a\mid x)
=\eta\pi_{g_t}(a\mid x)\left[(g'-g)(x,a)-\sum_{b\in\mathcal A}\pi_{g_t}(b\mid x)(g'-g)(x,b)\right].
\]
Therefore
\[
\sum_{a\in\mathcal A}\left|\frac{d}{dt}\pi_{g_t}(a\mid x)\right|\le 2\eta v,
\]
and integration over \(t\in[0,1]\) gives
\[
\lVert\pi_g(\cdot\mid x)-\pi_{g'}(\cdot\mid x)\rVert_1\le2\eta v.
\]
If also \(\lVert f-f'\rVert_\infty\le v\), then \(g-f\) and \(g'-f'\) take values in \([-1,1]\), whence
\[
\left|(g-f)^2-(g'-f')^2\right|
\le2\left|(g-f)-(g'-f')\right|\le4v
\]
pointwise. Splitting the difference between the Gibbs weights and the squared contrasts consequently shows
\[
\lVert H_{g,f}-H_{g',f'}\rVert_\infty\le(2\eta+4)v.
\]
Choose \(v=\{n(2\eta+4)\}^{-1}\). The Cartesian square of a \(v\)-net of \(\mathcal F\) induces a \(1/n\)-net for \(\{H_{g,f}:g,f\in\mathcal F\}\) whose cardinality \(M\) satisfies
\[
M\le\{1+4n(\eta+2)\}^{2d}.
\]

It remains to prove and apply a fixed-function deviation inequality. Let \(W_1,\ldots,W_n\) be independent random variables in \([0,1]\) with common mean \(\mu\). For \(|\lambda|<3\) and \(|w|\le1\), expansion of the exponential and the inequality \(k!\ge2\cdot3^{k-2}\) for \(k\ge2\) give
\[
e^{\lambda w}\le1+\lambda w+\frac{\lambda^2w^2}{2(1-|\lambda|/3)}.
\]
Apply this inequality to \(w=W_i-\mu\). Since \(\mathbb E(W_i-\mu)=0\) and \(\mathbb E(W_i-\mu)^2\le\mu\), independence and \(1+s\le e^s\) imply the corresponding Bernstein moment-generating-function bounds for both signs. Exponential Markov's inequality, with the optimizing parameter for the sub-gamma variance factor \(n\mu\) and scale \(1/3\), then gives, for every \(t>0\),
\[
\Pr\left\{\left|P_nW-\mu\right|>\sqrt{\frac{2\mu t}{n}}+\frac{t}{3n}\right\}\le2e^{-t}.
\]
The elementary inequality \(\sqrt{2\mu t/n}\le\mu/2+t/n\) shows that, outside this event,
\[
\mu\le2P_nW+\frac{8t}{3n},
\qquad
P_nW\le2\mu+\frac{8t}{3n}.
\]

Apply this two-sided event to each of the \(M\) net functions and set
\[
t=\log\frac{2M}{\zeta}.
\]
For each fixed net function the single event containing both comparisons fails with probability at most \(2e^{-t}\). Hence the union bound gives all comparisons on the net simultaneously with probability at least \(1-2Me^{-t}=1-\zeta\). For an arbitrary \(H_{g,f}\), choose a net function \(H_0\) with \(\lVert H_{g,f}-H_0\rVert_\infty\le1/n\). Replacing both its population and empirical means by those of \(H_{g,f}\) costs at most \(3/n\) in either comparison. Finally,
\[
t\le\log\frac2\zeta+2d\log\{1+4n(\eta+2)\}.
\]
Thus, on the same event and simultaneously for every \(g,f\in\mathcal F\),
\[
P_\rho H_{g,f}\le2P_nH_{g,f}+\frac{8}{3n}\left[\log\frac2\zeta+2d\log\{1+4n(\eta+2)\}\right]+\frac3n,
\]
\[
P_nH_{g,f}\le2P_\rho H_{g,f}+\frac{8}{3n}\left[\log\frac2\zeta+2d\log\{1+4n(\eta+2)\}\right]+\frac3n.
\]
This is the proposed sharpened statement.